\section{统计稳定性：从数据到概念的判别准则}\label{sec:statistical-stability}

本节沿用第 \ref{sec:preliminaries} 节（\ref{subsec:freq-setting}--\ref{def:stable-concept}）所固定的保测读出口径与稳定概念，并在该口径上给出用于判别与审计的分辨率指标、误差分解与可复现实验协议。

\subsection{稳定性与有限分辨率}

有限 $m$ 下微态空间 $\Omega_m$ 过大时，直接以原始微态为概念将导致不可控的稀疏性与不稳定。为建立可递归的概念层级，需要一个将 $\Omega_m$ 折叠到更小稳定类型空间的机制（第 \ref{sec:folding} 节）。该机制的数学要求是：折叠后类型空间随 $m$ 增长呈现自相似递归，并与黄金比例一致。

\subsection{定量拟合：稳定类型分布的偏差指标}\label{sec:quant-fit}
给定一段观测序列 $\omega_m(x;t)\in\Omega_m$，折叠后得到稳定类型序列（见第 \ref{sec:folding} 节）
\[
w_t:=\Fold_m(\omega_m(x;t))\in X_m.
\]
令 $\hat{\pi}_m$ 为经验分布（直方图）
\[
\hat{\pi}_m(u):=\frac{1}{N}\sum_{t=0}^{N-1}\ind_{\{w_t=u\}},\qquad u\in X_m.
\]
另一方面，黄金均值约束 $X_\infty$ 的最大熵（Parry）平衡态由 Perron--Frobenius 本征结构给出（见第 \ref{sec:group-unification} 节及注 \ref{prop:parry-measure-golden}；亦参见 \cite{LindMarcus1995}），其在长度 $m$ 上诱导的柱分布记为 $\pi_m$。于是可用如下两个无量纲偏差作为拟合偏差审计量：
\[
D_{\mathrm{TV}}(\hat{\pi}_m,\pi_m):=\frac12\sum_{u\in X_m}\abs{\hat{\pi}_m(u)-\pi_m(u)},\qquad
D_{\mathrm{KL}}(\hat{\pi}_m\|\pi_m):=\sum_{u\in X_m}\hat{\pi}_m(u)\log\frac{\hat{\pi}_m(u)}{\pi_m(u)}.
\]
关于 Parry 分布的构造与标准性质见 \cite{LindMarcus1995}。

\subsection{源模型与折叠输出的记号（仅用于定位对象）}\label{subsec:histogram-convergence-natural-sources}
本节固定“源 $\Rightarrow$ 折叠输出直方图”的记号，并记录后文直接调用的标准收敛/置信半径形式；推导见 \cite{WaltersErgodicTheory,WeissmanEtAl2003L1Empirical}。

\begin{definition}[一般源与折叠输出过程]\label{def:source-and-fold-output}
令 $S=(S_t)_{t\in\ZZ}$ 为取值于 $\{0,1\}$ 的双边随机过程，定义长度-$m$ 窗口与折叠输出为
\[
W_t^{(m)}:=S_tS_{t+1}\cdots S_{t+m-1}\in\Omega_m,
\qquad
Y_t^{(m)}:=\Fold_m(W_t^{(m)})\in X_m.
\]
对样本长度 $N$，定义折叠直方图（经验分布）
\[
\hat{\pi}_{m,N}(u):=\frac{1}{N}\sum_{t=0}^{N-1}\ind_{\{Y_t^{(m)}=u\}},
\qquad u\in X_m,
\]
并记真实边缘分布（极限目标）为
\[
\pi_m^{\mathrm{fold}}(u):=\mathbb{P}\bigl(Y_0^{(m)}=u\bigr),\qquad u\in X_m.
\]
\end{definition}

若 $S$ 平稳且遍历，则对任意固定 $m$ 与任意 $u\in X_m$ 有
\[
\hat{\pi}_{m,N}(u)\ \xrightarrow[N\to\infty]{}\ \pi_m^{\mathrm{fold}}(u)
\quad\text{几乎处处}
\]
（\cite{WaltersErgodicTheory}）。

设 $S_t$ 为 i.i.d. 伯努利源。对重叠窗口的折叠直方图，其有效样本量降为 $N_{\min}=\lfloor N/m\rfloor$，并可由 \cite{WeissmanEtAl2003L1Empirical} 的 TV 置信半径给出显式指数型证书：
\[
\varepsilon_{\mathrm{TV}}(N_{\min},|X_m|;\delta)=\frac12\sqrt{\frac{2}{N_{\min}}\Bigl(|X_m|\log 2+\log(1/\delta)\Bigr)}.
\]

\begin{remark}[与 IID 块采样实验的关系]\label{rem:iid-overlap-vs-blocks}
第 \ref{sec:experiments} 节的 IID 基线实验采用独立 $m$-块采样（有效样本量为 $N$）；沿时间滑窗采样时有效样本量为 $N_{\min}=\lfloor N/m\rfloor$。
\end{remark}

\subsection{有限分辨率同余刚性与筛分障碍}\label{subsec:stat-stability-finite-resolution-congruence-obstructions}

本小节把第 \ref{sec:folding} 节所建立的模余类刚性（推论 \ref{cor:fold-mod-rigidity-periodic-predicates}）在“判别/筛分”语境下改写为两类直接后果：一类是周期 Dirichlet 级数的 Hurwitz \(\zeta\) 分解，另一类是固定分辨率余类筛在素数判别中的精确率退化。

\begin{proposition}[稳定谓词的余类化与周期系数]\label{prop:stable-predicate-residue-classification}
固定 \(m\ge 1\)，令 \(M:=F_{m+2}\)。
对任意谓词 \(P\subseteq X_m\)，定义其余类像
\[
R_P\ :=\ \{\,v_m(x)\bmod M:\ x\in P\,\}\ \subseteq\ \ZZ/M\ZZ,
\]
并定义周期系数
\[
a_P(n)\ :=\ \ind_{\{n\bmod M\in R_P\}}\qquad(n\ge 1).
\]
则对任意 \(\omega\in\Omega_m\) 有
\[
\ind_{\{\Fold_m(\omega)\in P\}}\ =\ a_P\!\bigl(N(\omega)\bigr),
\]
从而 \(a_P\) 为周期 \(M\) 的算术函数。
\end{proposition}

\begin{proof}
由推论 \ref{cor:fold-mod-rigidity-periodic-predicates} 的因子化 \(\Fold_m=s_m\circ r_m\)，取 \(R_P:=s_m^{-1}(P)\) 即得
\(
\Fold_m(\omega)\in P\iff r_m(\omega)\in R_P
\)，
而 \(r_m(\omega)=N(\omega)\bmod M\)。
将其写成指示函数即得结论。证毕。
\end{proof}

\begin{theorem}[周期 Dirichlet 级数的 Hurwitz \texorpdfstring{$\zeta$}{zeta} 分解]\label{thm:periodic-dirichlet-series-hurwitz-decomposition}
在命题 \ref{prop:stable-predicate-residue-classification} 的记号下，对 \(\Re(s)>1\) 定义 Dirichlet 级数
\[
D_P(s)\ :=\ \sum_{n\ge 1}\frac{a_P(n)}{n^s}.
\]
则成立显式分解
\[
D_P(s)\ =\ M^{-s}\sum_{r=1}^{M} a_P(r)\,\zeta\!\left(s,\frac{r}{M}\right),
\]
其中 \(\zeta(s,\alpha)\) 为 Hurwitz \(\zeta\) 函数（\(\alpha\in(0,1]\)）。
因此 \(D_P\) 的解析延拓与奇点结构由有限个 Hurwitz \(\zeta\) 的线性组合完全决定（参见 \cite{Apostol1976,Titchmarsh1986RiemannZeta}）。
\end{theorem}

\begin{proof}
由周期性分解 \(n=Mk+r\)（其中 \(k\ge 0\)，\(1\le r\le M\)），得到
\[
\sum_{n\ge 1}\frac{a_P(n)}{n^s}
=\sum_{r=1}^{M}a_P(r)\sum_{k\ge 0}\frac{1}{(Mk+r)^s}
=M^{-s}\sum_{r=1}^{M}a_P(r)\sum_{k\ge 0}\frac{1}{\left(k+\frac{r}{M}\right)^s}.
\]
内层和即为 \(\zeta(s,r/M)\)。证毕。
\end{proof}

\begin{remark}[Euler 积阻断的结构性原因]\label{rem:euler-product-obstruction-periodic-indicator}
若 \(D_P(s)\) 在 \(\Re(s)>1\) 上具有收敛的 Euler 积表示，则其系数 \(a_P(n)\) 必为（至少）乘法函数。
对一般余类指示函数 \(a_P(n)=\ind_{\{n\bmod M\in R_P\}}\)，该乘法性通常不成立，因此不会出现“跨所有素数的”Euler 积结构；只有在极少数退化情形（例如 \(a_P\equiv 0\)、\(a_P\equiv 1\)，或 \(a_P(n)=\ind_{\{\gcd(n,M)=1\}}\) 这类由主特征诱导的情形）才可能出现有限因子修正的 Euler 积。
\end{remark}

\begin{theorem}[固定分辨率余类筛的素数精确率退化]\label{thm:fixed-resolution-prime-sieve-precision-to-zero}
固定 \(m\ge 1\)，令 \(M:=F_{m+2}\)，并取任意 \(P\subseteq X_m\)。
令 \(R_P\subseteq\ZZ/M\ZZ\) 与 \(a_P\) 如命题 \ref{prop:stable-predicate-residue-classification}。
定义接受集合
\[
S_P\ :=\ \{n\in\NN:\ a_P(n)=1\}.
\]
若 \(R_P\) 含至少一个与 \(M\) 互素的余类，则对经验精确率
\[
\mathrm{Prec}_P(X)\ :=\ \frac{\#\{p\le X:\ p\ \text{为素数且}\ p\in S_P\}}{\#\{n\le X:\ n\in S_P\}}
\]
成立渐近上界
\[
\mathrm{Prec}_P(X)\ \le\ \frac{M}{\varphi(M)}\cdot \frac{1+o(1)}{\log X}\qquad(X\to\infty),
\]
从而 \(\mathrm{Prec}_P(X)\to 0\)。
\end{theorem}

\begin{proof}
记 \(U_M:=\{a\in\ZZ/M\ZZ:\gcd(a,M)=1\}\)。
由算术级数中的素数定理（参见 \cite{Titchmarsh1986RiemannZeta}），对每个 \(a\in R_P\cap U_M\) 有
\[
\#\{p\le X:\ p\equiv a\ (\mathrm{mod}\ M)\}
=\frac{1+o(1)}{\varphi(M)}\frac{X}{\log X}.
\]
求和得
\[
\#\{p\le X:\ p\in S_P\}
=\frac{\card{R_P\cap U_M}}{\varphi(M)}\frac{1+o(1)}{\log X}\,X.
\]
另一方面，余类计数给出
\[
\#\{n\le X:\ n\in S_P\}=\frac{\card{R_P}}{M}X+O(1).
\]
相除得到
\[
\mathrm{Prec}_P(X)
=\frac{\card{R_P\cap U_M}}{\card{R_P}}\cdot\frac{M}{\varphi(M)}\cdot\frac{1+o(1)}{\log X}
\ \le\ \frac{M}{\varphi(M)}\cdot\frac{1+o(1)}{\log X},
\]
从而趋于 \(0\)。证毕。
\end{proof}

\begin{conjecture}[素数集合的有限分辨率 lift--obstruction]\label{conj:prime-lift-obstruction-finite-resolution}
对每个 \(m\ge 1\)，令 \(\mathcal B_m\subseteq\mathcal P(\NN)\) 为由模 \(F_{m+2}\) 的余类集合生成的布尔代数（等价地，由所有 \(S_P\) 生成，其中 \(P\subseteq X_m\) 任意）。
记 \(\mathcal B_\infty:=\bigcup_{m\ge 1}\mathcal B_m\)。
猜想不存在一族集合 \(A_m\in\mathcal B_m\) 使得其在自然投影意义下相容，并且其指示函数在任意固定 \(s>1\) 的 Dirichlet 拓扑中逼近素数指示函数 \(\ind_{\PP}\)。
\end{conjecture}


